\subsubsection{Spin-Axis Projectors}
The left action of a chosen unit axis \(\vct{u}\in\R^3\) is resolved by the
projector family from Eq.\ \eqref{eq:projectors}. For clarity, that
definition is restated here as an ordinary equality:
\[
\pv{\Pi}(\pm)=\tfrac12(1\pm\vct{u}\cdot\sigv),
\qquad
\vct{u}^{\,2}=1,
\qquad
\text{from Eq.\ \eqref{eq:projectors}.}
\]
In addition to the
idempotence, orthogonality, completeness, and axis-eigenvalue identities
listed there, the following involution actions are obtained:
\begin{align*}
\pv{\Pi}(\pm)^{\mathsf H}&=\pv{\Pi}(\pm),\\
\pv{\Pi}(\pm)^*&=\pv{\Pi}(\mp).
\end{align*}
Accordingly,
\begin{equation*}
\vct{u}\cdot\sigv=(+1)\pv{\Pi}(+)+(-1)\pv{\Pi}(-),
\end{equation*}
so the eigenvalues \(+1\) and \(-1\) of the axis operator are carried by the
two projector parts.

In the Pauli matrix realization, each \(\pv{\Pi}(\pm)\) is of rank one. A general
element of \(\C+\C^3\), however, is a full \(2\times2\) complex matrix,
not one selected two-entry column. The projection formulas are therefore
valid column by column. A single Pauli wave-function column is described after a
fixed projector is placed on the right; without that restriction, two
independent columns are contained in the full paravector.

If the axis is rotated by the rotor \(\pv{R}\) from Eq.\ \eqref{eq:rotor},
the transformed axis and its second projector pair satisfy
\begin{equation}
\vct{u}'\cdot\sigv=\herm{\pv{R}}(\vct{u}\cdot\sigv)\pv{R},
\qquad
\pv{\Pi}'(\pm)=\tfrac12(1\pm \vct{u}'\cdot\sigv)
=\herm{\pv{R}}\pv{\Pi}(\pm)\pv{R}.
\label{eq:rotated-axis-projectors}
\end{equation}
Then
\begin{equation*}
\vct{u}'\cdot\sigv=(+1)\pv{\Pi}'(+)+(-1)\pv{\Pi}'(-).
\end{equation*}

For an algebra element \(\pv{\psi}\in\C+\C^3\), specialize the left-projected
parts from Eq.\ \eqref{eq:projectors}:
\begin{equation*}
\pv{\psi}(\pm)=\pv{\Pi}(\pm)\pv{\psi}.
\end{equation*}
The following identities are obtained:
\begin{align*}
(\vct{u}\cdot\sigv)\pv{\psi}(\pm)&=\pm\pv{\psi}(\pm),\\
\pv{\psi}&=\pv{\psi}(+)+\pv{\psi}(-).
\end{align*}
Thus each spinor column is split into the two eigenspaces of the chosen axis.
The Hermitian operation is represented by
\begin{align*}
\herm{\pv{\psi}(\pm)}
&=\herm{\pv{\psi}}\,\pv{\Pi}(\pm)^{\mathsf H}\notag\\
&=\herm{\pv{\psi}}\,\pv{\Pi}(\pm),
\end{align*}
and therefore
\begin{align*}
\herm{\pv{\psi}(\pm)}\pv{\psi}(\pm)&=\herm{\pv{\psi}}\,\pv{\Pi}(\pm)\pv{\psi},\\
\herm{\pv{\psi}(\pm)}\pv{\psi}(\mp)&=\herm{\pv{\psi}}\,\pv{\Pi}(\pm)\pv{\Pi}(\mp)\pv{\psi}=0,\\
\herm{\pv{\psi}(\mp)}\pv{\psi}(\pm)&=\herm{\pv{\psi}}\,\pv{\Pi}(\mp)\pv{\Pi}(\pm)\pv{\psi}=0.
\end{align*}
For two algebra elements \(\pv{\phi},\pv{\psi}\), the pointwise algebraic
pairing associated with \(\pv{M}\) is
\(\eqtext{scalar}(\herm{\pv{\phi}}\pv{M}\pv{\psi})\). For a nonzero
\(\pv{\psi}\), the corresponding normalized diagonal matrix element is
\begin{equation*}
\frac{\eqtext{scalar}(\herm{\pv{\psi}}\pv{M}\pv{\psi})}
     {\eqtext{scalar}(\herm{\pv{\psi}}\pv{\psi})}.
\end{equation*}
When \(\pv{M}\) is Hermitian for the relevant pairing, this
diagonal element is real and is an expectation value. For a general
\(\pv{M}\), it is only a normalized matrix element and may be complex.
Scalar extraction is one half of the matrix trace. Consequently, for a full
matrix \(\pv{\psi}\) the numerator and denominator are each one half of the
sum of their two ordinary column pairings. After one column is selected by a
fixed right projector, they are each one half of the corresponding ordinary
selected-column quantity. The common factor is canceled in the normalized
matrix element. The spatially integrated quantum bracket is defined
separately in Eq.\ \eqref{eq:adjoint-paravector-braket}.

Rotational covariance is made explicit by the same convention. Under the
active rotation
\begin{equation*}
\pv{\psi}'=\herm{\pv{R}}\pv{\psi},
\qquad
\pv{M}'=\herm{\pv{R}}\pv{M}\pv{R},
\qquad
\pv{R}\herm{\pv{R}}=1,
\end{equation*}
both the numerator and denominator of the normalized matrix element are
unchanged:
\begin{equation}
\frac{\eqtext{scalar}(\herm{\pv{\psi}'}\pv{M}'\pv{\psi}')}
     {\eqtext{scalar}(\herm{\pv{\psi}'}\pv{\psi}')}
=
\frac{\eqtext{scalar}(\herm{\pv{\psi}}\pv{M}\pv{\psi})}
     {\eqtext{scalar}(\herm{\pv{\psi}}\pv{\psi})}.
\label{eq:rotor-expectation-covariance}
\end{equation}
Equivalently, if the observable \(\pv{M}\) is held fixed while only the
carrier is actively rotated, its value in the original carrier is computed
with the rotated-back observable \(\pv{R}\pv{M}\herm{\pv{R}}\). This is the
active-versus-passive equivalence expressed in
Eq.\ \eqref{eq:rotated-axis-projectors}.

For an axis observable, the following identity is obtained by cyclic scalar
extraction:
\begin{equation}
\eqtext{scalar}\!\left(
  \herm{\pv{\psi}}(\vct{u}\cdot\sigv)\pv{\psi}
\right)
=
\vct{u}\cdot
\eqtext{vector}\!\left(\pv{\psi}\herm{\pv{\psi}}\right).
\label{eq:axis-polarization-identity}
\end{equation}
Consequently its normalized pointwise value is
\begin{equation*}
\frac{\vct{u}\cdot
  \eqtext{vector}(\pv{\psi}\herm{\pv{\psi}})}
 {\eqtext{scalar}(\herm{\pv{\psi}}\pv{\psi})}.
\end{equation*}
In the selected-column matrix realization, twice the numerator is the
ordinary column matrix element and twice the denominator is the ordinary
column norm, so the common factor is canceled. For a spatial wave field the
same statement is valid with both quantities integrated as in
Eqs.\ \eqref{eq:adjoint-paravector-braket} and \eqref{eq:expectation-map}.

The unnormalized pointwise numerator for the rotated axis is expanded as
\begin{align*}
\eqtext{scalar}\!\left(\herm{\pv{\psi}}(\vct{u}'\cdot\sigv)\pv{\psi}\right)
&=\eqtext{scalar}\!\left(\herm{\pv{\psi}(+)}(\vct{u}'\cdot\sigv)\pv{\psi}(+)\right)\notag\\
&\quad+\eqtext{scalar}\!\left(\herm{\pv{\psi}(+)}(\vct{u}'\cdot\sigv)\pv{\psi}(-)\right)\notag\\
&\quad+\eqtext{scalar}\!\left(\herm{\pv{\psi}(-)}(\vct{u}'\cdot\sigv)\pv{\psi}(+)\right)\notag\\
&\quad+\eqtext{scalar}\!\left(\herm{\pv{\psi}(-)}(\vct{u}'\cdot\sigv)\pv{\psi}(-)\right).
\end{align*}

The algebra governing these same-side and cross-side terms is
\begin{align*}
(\vct{u}'\cdot\sigv)(\vct{u}\cdot\sigv)
&=\vct{u}'\cdot\vct{u}+i(\vct{u}'\times\vct{u})\cdot\sigv,\\
(\vct{u}\cdot\sigv)(\vct{u}'\cdot\sigv)
&=\vct{u}'\cdot\vct{u}-i(\vct{u}'\times\vct{u})\cdot\sigv,\\
(\vct{u}\cdot\sigv)(\vct{u}'\cdot\sigv)(\vct{u}\cdot\sigv)
&=\bigl(2(\vct{u}'\cdot\vct{u})\vct{u}-\vct{u}'\bigr)\cdot\sigv.
\end{align*}
The second axis is decomposed into its \(\eqtext{par}\) and
\(\eqtext{prp}\) parts by the axis functions from
Eqs.\ \eqref{eq:par-function} and \eqref{eq:prp-function}:
\begin{equation*}
\vct{u}'=\eqtext{par}(\vct{u}')+\eqtext{prp}(\vct{u}'),
\qquad
\eqtext{par}(\vct{u}')=(\vct{u}\cdot\vct{u}')\vct{u}.
\end{equation*}
Then the projector products are
\begin{equation*}
\pv{\Pi}(\pm)(\vct{u}'\cdot\sigv)\pv{\Pi}(\pm)=\pm(\vct{u}\cdot\vct{u}')\pv{\Pi}(\pm),
\end{equation*}
\begin{equation*}
\pv{\Pi}(\pm)(\vct{u}'\cdot\sigv)\pv{\Pi}(\mp)
=\frac12\Bigl(\eqtext{prp}(\vct{u}')\pm i\,\vct{u}\times\vct{u}'\Bigr)\cdot\sigv\,\pv{\Pi}(\mp).
\end{equation*}
The diagonal terms are therefore weighted by the axis overlap \(\vct{u}\cdot\vct{u}'\), while the off-diagonal terms are controlled by the \(\eqtext{prp}\) and commutator parts of the rotated axis.

If \(\theta\) is the angle between \(\vct{u}\) and \(\vct{u}'\), then
\begin{equation*}
\vct{u}'\cdot\vct{u}=\cos(\theta),
\end{equation*}
Because scalar extraction is one half of the matrix trace, the
trace-normalized rank-one projector overlap is
\begin{equation*}
2\,\eqtext{scalar}\!\bigl(\pv{\Pi}(s)\pv{\Pi}'(s')\bigr)
=\frac12\bigl(1+ss'\,\vct{u}\cdot\vct{u}'\bigr),
\qquad s,s'\in\{+1,-1\}.
\end{equation*}
For a normalized irreducible axis eigenstate this overlap is the transition
probability. In particular, the same-side and opposite-side probabilities
are
\begin{equation*}
\frac12(1+\vct{u}'\cdot\vct{u})=\cos^2(\theta/2),
\qquad
\frac12(1-\vct{u}'\cdot\vct{u})=\sin^2(\theta/2).
\end{equation*}
If \(\vct{u}'\) is obtained by rotating \(\vct{u}\) through angle \(\phi\)
about a unit axis \(\vct{n}\), the overlap is given by Rodrigues's formula:
\begin{align*}
\vct{u}'\cdot\vct{u}
&=(\vct{u}\cdot\vct{u})\cos(\phi)
+\bigl(\vct{u}\cdot(\vct{n}\times\vct{u})\bigr)\sin(\phi)\notag\\
&\quad+(\vct{n}\cdot\vct{u})\bigl(1-\cos(\phi)\bigr)(\vct{n}\cdot\vct{u})\notag\\
&=(\vct{n}\cdot\vct{u})^2+\bigl(1-(\vct{n}\cdot\vct{u})^2\bigr)\cos(\phi).
\end{align*}
In the perpendicular case \(\vct{n}\cdot\vct{u}=0\), this is reduced to
\begin{equation*}
\vct{u}'\cdot\vct{u}=\cos(\phi),
\end{equation*}
hence
\begin{equation*}
\frac12(1+\vct{u}'\cdot\vct{u})=\cos^2(\phi/2),
\qquad
\frac12(1-\vct{u}'\cdot\vct{u})=\sin^2(\phi/2).
\end{equation*}
These are the familiar spin-\(\tfrac12\) transition probabilities. Without
normalization and the irreducible-state qualification, only projector
overlaps are represented.

\subsubsection{Ladder Operations About the z Axis}
For raising and lowering about the \(z\) axis, the following axis and
projector pair are chosen:
\begin{equation*}
\vct{u}=\vct{e}_3=(0,0,1),
\qquad
\pv{\Pi}(\pm)=\tfrac12(1\pm \pv{\sigma}_3).
\end{equation*}
The ladder elements already defined in Eq.\ \eqref{eq:ladder} are used:
\begin{equation*}
\pv{a}(\pm)=\tfrac12(\pv{\sigma}_1\pm i\pv{\sigma}_2).
\end{equation*}
The commutator relation is satisfied:
\begin{equation*}
\com{\pv{\sigma}_3}{\pv{a}(\pm)}=\pm 2\pv{a}(\pm),
\end{equation*}
so the \(\pv{\sigma}_3\) eigenvalue is raised by \(\pv{a}(+)\) and lowered
by \(\pv{a}(-)\).

Their projector products are
\begin{equation*}
\pv{a}(+)\pv{a}(-)=\pv{\Pi}(+),
\qquad
\pv{a}(-)\pv{a}(+)=\pv{\Pi}(-),
\end{equation*}
and each ladder operator is nilpotent:
\begin{equation*}
\pv{a}(\pm)^2=0.
\end{equation*}
The one-sided projector actions are
\begin{equation*}
\begin{aligned}
\pv{a}(+)\pv{\Pi}(+) &= 0, &
\pv{a}(+)\pv{\Pi}(-) &= \pv{a}(+), \\
\pv{\Pi}(+)\pv{a}(+) &= \pv{a}(+), &
\pv{\Pi}(-)\pv{a}(+) &= 0, \\
\pv{a}(-)\pv{\Pi}(-) &= 0, &
\pv{a}(-)\pv{\Pi}(+) &= \pv{a}(-), \\
\pv{\Pi}(-)\pv{a}(-) &= \pv{a}(-), &
\pv{\Pi}(+)\pv{a}(-) &= 0,
\end{aligned}
\end{equation*}
by which the raising/lowering interpretation is made explicit. The original
sigma generators are recovered from the same ladder/projector data:
\begin{equation*}
\begin{aligned}
\pv{\sigma}_1&=\pv{a}(+)+\pv{a}(-),\qquad
\pv{\sigma}_2=\frac{\pv{a}(+)-\pv{a}(-)}{i},\\
\pv{\sigma}_3&=\pv{\Pi}(+)-\pv{\Pi}(-).
\end{aligned}
\end{equation*}

\subsubsection{Spin-Axis Noncommutativity}
For unit directions \(\vct{u},\vct{u}'\in\R^3\),
\begin{align*}
\acom{\vct{u}\cdot\sigv}{\vct{u}'\cdot\sigv}
&=2\,\vct{u}\cdot\vct{u}',\\
\com{\vct{u}\cdot\sigv}{\vct{u}'\cdot\sigv}
&=2i(\vct{u}\times\vct{u}')\cdot\sigv.
\end{align*}
The axis overlap is recorded by the anticommutator, while the oriented area
spanned by the two axes is recorded by the commutator. Thus exact commutation
of the axis operators is obtained when the directions are parallel or
antiparallel. These identities are purely algebraic. A normalized state,
expectation functional, and state-dependent variances are additionally
required for a Robertson or Robertson--Schr\"odinger uncertainty relation;
no such inequality is inferred from the
commutator alone \cite{robertson1929}.
